%%%%%%%%%%%%%%%%%%%%%%%%%%%%%%%%%%%%%%%%%
% Arsclassica Article
% Structure Specification File
%
% This file has been downloaded from:
% http://www.LaTeXTemplates.com
%
% Original author:
% Lorenzo Pantieri (http://www.lorenzopantieri.net) with extensive modifications by:
% Vel (vel@latextemplates.com)
%
% License:
% CC BY-NC-SA 3.0 (http://creativecommons.org/licenses/by-nc-sa/3.0/)
%
%%%%%%%%%%%%%%%%%%%%%%%%%%%%%%%%%%%%%%%%%

%----------------------------------------------------------------------------------------
%	REQUIRED PACKAGES
%----------------------------------------------------------------------------------------

\usepackage[
nochapters, % Turn off chapters since this is an article        
beramono, % Use the Bera Mono font for monospaced text (\texttt)
%eulermath,% Use the Euler font for mathematics
pdfspacing, % Makes use of pdftex’ letter spacing capabilities via the microtype package
dottedtoc % Dotted lines leading to the page numbers in the table of contents
]{classicthesis} % The layout is based on the Classic Thesis style

\usepackage{arsclassica} % Modifies the Classic Thesis package

\usepackage[T1]{fontenc} % Use 8-bit encoding that has 256 glyphs

\usepackage[utf8]{inputenc} % Required for including letters with accents

\usepackage{graphicx} % Required for including images
\graphicspath{{Figures/}} % Set the default folder for images

\usepackage{enumitem} % Required for manipulating the whitespace between and within lists

\usepackage{lipsum} % Used for inserting dummy 'Lorem ipsum' text into the template

% \usepackage{subfig} % Required for creating figures with multiple parts (subfigures)

\usepackage{amsmath,amssymb,amsthm} % For including math equations, theorems, symbols, etc

\usepackage{varioref} % More descriptive referencing




%----------------------------------------------------------------------------------------
% MY PACKAGES
%----------------------------------------------------------------------------------------
\usepackage{microtype}
%\usepackage{tgheros}
%%% uni formatting
%\renewcommand{\familydefault}{\sfdefault} % when reverting back add eulermath to classicthesis option
\usepackage{newtxsf}
\linespread{1.15}
\usepackage{lineno}


% code to stop section line numbering from substack
% Source - https://tex.stackexchange.com/a/635606
% Posted by Jinwen
% Retrieved 2026-09-03, License - CC BY-SA 4.0
\newif\ifLNturnsON
\def\LocallyStopLineNumbers{\LNturnsONfalse%
	\ifLineNumbers\LNturnsONtrue\fi\nolinenumbers}
\def\ResumeLineNumbers{\ifLNturnsON\linenumbers\fi}

\makeatletter

\titleformat{\section}
{\LocallyStopLineNumbers\normalfont\Large\bfseries}{\thesection}{1em}{}[\ResumeLineNumbers]
\titleformat{\subsection}
{\LocallyStopLineNumbers\normalfont\large\bfseries}{\thesubsection}{1em}{}[\ResumeLineNumbers]
\titleformat{\subsubsection}
{\LocallyStopLineNumbers\normalfont\normalsize\bfseries}{\thesubsubsection}{1em}{}[\ResumeLineNumbers]
\titleformat{\paragraph}[runin]
{\LocallyStopLineNumbers\normalfont\normalsize\bfseries}{\theparagraph}{1em}{}[\ResumeLineNumbers]
\titleformat{\subparagraph}[runin]
{\LocallyStopLineNumbers\normalfont\normalsize\bfseries}{\thesubparagraph}{1em}{}[\ResumeLineNumbers]
%%% end of uni formatting

\usepackage[mathlf]{MinionPro} % onlytext
\usepackage{float}
%\usepackage[margin=3cm]{geometry}   % for non-printing
\usepackage[twoside,inner=3.5cm,outer=2.5cm,top=3cm,bottom=3cm]{geometry}
\usepackage{subcaption}
\usepackage{caption}
\captionsetup{
	format=plain,      % don't use the hanging/justified style
	justification=justified,
	singlelinecheck=false,
	labelsep=colon,    % "Figure 1: " inline before the text
	labelfont=bf,      % optional — bold "Figure 1:"
	textfont=normalfont
}

\usepackage[style=authoryear,
 maxbibnames=10,
 minbibnames=10,
 maxcitenames=2]{biblatex}
 
 % italicised 'et al.' from stackexchange:
 \renewbibmacro*{name:andothers}{% Based on name:andothers from biblatex.def
 	\ifboolexpr{
 		test {\ifnumequal{\value{listcount}}{\value{liststop}}}
 		and
 		test \ifmorenames
 	}
 	{\ifnumgreater{\value{liststop}}{1}
 		{\finalandcomma}
 		{}%
 		\andothersdelim\bibstring[\emph]{andothers}}
 	{}}
 	
 	\setlength\bibitemsep{0.4\baselineskip} % spacing in reference list
 
\addbibresource{references.bib}
\usepackage{siunitx}
\usepackage{gensymb}
\usepackage{nicefrac}
\usepackage{tabularx}
\usepackage[table]{xcolor}
\usepackage{acro}
\usepackage{tikz}





%----------------------------------------------------------------------------------------
%	THEOREM STYLES
%---------------------------------------------------------------------------------------

\theoremstyle{definition} % Define theorem styles here based on the definition style (used for definitions and examples)
\newtheorem{definition}{Definition}

\theoremstyle{plain} % Define theorem styles here based on the plain style (used for theorems, lemmas, propositions)
\newtheorem{theorem}{Theorem}

\theoremstyle{remark} % Define theorem styles here based on the remark style (used for remarks and notes)

%----------------------------------------------------------------------------------------
%	HYPERLINKS
%---------------------------------------------------------------------------------------

\hypersetup{
%draft, % Uncomment to remove all links (useful for printing in black and white)
colorlinks=true, breaklinks=true, bookmarks=true,bookmarksnumbered,
urlcolor=webbrown, linkcolor=RoyalBlue, citecolor=webgreen, % Link colors
pdftitle={}, % PDF title
pdfauthor={\textcopyright}, % PDF Author
pdfsubject={}, % PDF Subject
pdfkeywords={}, % PDF Keywords
pdfcreator={pdfLaTeX}, % PDF Creator
pdfproducer={LaTeX with hyperref and ClassicThesis} % PDF producer
}



%----------------------------------------------------------------------------------------
%	ACRONYMS
%---------------------------------------------------------------------------------------




\DeclareAcronym{ET}{
	short=ET,
	long=Evapotranspiration,
}
\DeclareAcronym{SHAP}{
	short=SHAP,
	long=SHapley Additive exPlanations,
}
\DeclareAcronym{ML}{
	short=ML,
	long=Machine Learning,
}
\DeclareAcronym{RF}{
	short=RF,
	long=Random Forest,
}
\DeclareAcronym{DL}{
	short=DL,
	long=Deep Learning,
}
\DeclareAcronym{EC}{
	short=EC,
	long=Eddy Covariance,
}
\DeclareAcronym{PAM}{
	short=PAM,
	long=Partitioning Around Medoids,
}
\DeclareAcronym{R2}{
	short={$R^2$},
	long=Coefficient of determination,
}
\DeclareAcronym{RMSE}{
	short=RMSE,
	long=Root Mean Square Error,
}
\DeclareAcronym{MAE}{
	short=MAE,
	long=Mean Absolute Error,
}
\DeclareAcronym{LOGO}{
	short=LOGO,
	long=Leave One Group Out,
}
\DeclareAcronym{NEP}{
	short=NEP,
	long=Net Ecosystem Productivity,
}
\DeclareAcronym{LAI}{
	short=LAI,
	long=Leaf Area Index,
}
\DeclareAcronym{RNN}{
	short=RNN,
	long=Recurrent Neural Network,
}
\DeclareAcronym{ANN}{
	short=ANN,
	long=Artificial Neural Network,
}
\DeclareAcronym{SVM}{
	short=SVM,
	long=Support Vector Machine,
}
\DeclareAcronym{RS}{
	short=RS,
	long=Remote Sensing,
}
\DeclareAcronym{VPD}{
	short=VPD,
	long=Vapour Pressure Deficit,
}
\DeclareAcronym{SC}{
	short={$g_s$},
	long=Stomatal Conductance,
}


